\documentclass[UTF8]{ctexart}
\usepackage{graphicx}
\usepackage{float}
\usepackage{xcolor}

\title{TeXButler 中文样例}
\author{TeXButler 团队}
\date{\today}

\begin{document}

\maketitle

\section{引言}

这是一个用于回归测试的中文 ctex 文档。
它包含中文 LaTeX 特有的常见踩坑点（见下面各节），供规则引擎自测。

\section{常见错误示例}

\subsection{百分号误当注释}

本次测试通过率 71%，同比增长 5%。

正确写法应为 71\% 与 5\%。

\subsection{斜体包中文}

\textit{这是中文斜体，中文字体通常没有斜体}，应改用 \textbf{加粗} 或 {\bfseries 粗体}。

\subsection{浮动体定位}

\begin{figure}[ht]
\centering
\caption{使用 [ht] 的浮动体会乱飘}
\end{figure}

\begin{table}[H]
\centering
\caption{使用 [H] 的浮动体固定在本位}
\end{table}

\subsection{混色语法}

{\color{red!60} 需要 xcolor 的混色语法}。

\subsection{浮点垃圾}

正确率 87.30000000000001%，这是程序生成时未 round 的典型例子。

\section{结论}

本样例用于验证 TeXButler 的编译与规则检查功能。

\end{document}
